\scanpage{191}
% manuscript position: top
\mspage{107}
% manuscript position: 1
$\pi$-\unsure{accessibles|accessible}

% manuscript position: 2–4
Dans la suite, nous allons, plus généralement, considérer une catégorie $C$ équivalente à une petite catégorie, \inserted{stable par $\varinjlim$ de cardinal $\leq\pi_0$,\unsure{$^{\ast}$}} et la catégorie
% manuscript position: left/2–5
% slanted
\margin{il n'y a pas lieu de le \unsure{supp.} \unsure{ça} tout de suite}
% manuscript position: 5
\[
M\overset{\mathrm{déf}}{=}\operatorname{Ind}'_{\pi_0}C\underset{\unsure{\neq|\not\simeq}}{\subset}\widehat C,\tag*{(4.5.1)}
\]
% manuscript position: 6
où $\pi_0$ est un cardinal infini donné.

% manuscript position: 7–8
Je me propose d'introduire sur $M$ une filtration cardinale canonique.

% manuscript position: 9–11
\uline{Proposition} 4.5.2. Soit $F\in\operatorname{Ob}M\subset\operatorname{Ob}\widehat C$, et soit $\pi$ un cardinal $\geq\pi_1=\operatorname{Sup}(\pi_0,\operatorname{card}\operatorname{Fl} C_0)$, \inserted{où $C_0$ est une \inserted{petite} sous-catégorie pleine \deleted{\unsure{réduite}} de $C$ telle que $C_0\overset{\approx}{\hookrightarrow}C$.} Les conditions suivantes sont équivalentes
% manuscript position: left/9–13
% slanted
\margin{N'est \uline{prouvé} (et n'est vrai) que si $\pi^{\pi_0}=\pi$ \unsure{4.5.5|4.5.6}. Voir rectification page 116}

% manuscript position: 12
\[
\text{a)}\quad\operatorname{Card}\operatorname{Fl}(C_{0/F})\leq\pi
\]

% manuscript position: 13
\[
\text{a$'$)}\quad\operatorname{Card}\deleted{\unsure{F}}\operatorname{Ob}(C_{0/F}^{\unsure{\circ|)}})\leq\pi
\]

% manuscript position: 14
a$''$) Pour tout $x$ dans $C_0$, $\operatorname{card} F(x)\leq\pi$

% manuscript position: 15–18
b) Il existe une petite catégorie filtrante $I$ (\unsure{par une|par ex.} ordonnée), \uline{grande devant $\pi_0$}, \inserted{avec \uline{$\operatorname{card} I\leq\pi$},} et un système inductif $(x_i)_{i\in I}$ dans $C$, tel que l'on ait
% manuscript position: 19
\[
\deleted{\unsure{x}}F\simeq\varinjlim_{i\in I}x_i,\tag*{(4.5.2.1)}
\]
% manuscript position: 20–24
\deleted{\unsure{dans}} la $\varinjlim$ est prise dans $M$, ou ce qui revient au \unsure{$=$} (en vertu de 4.4.15, $I$ étant gde devant $\pi_0$) dans $\widehat C$, i.e. $F$ est isomorphe au foncteur ind-représentable défini par le Ind-objet $\underline{x}=(x_i)_{i\in I}$ de $C$.

% manuscript position: 25
c) $F$ est $\pi$-accessible dans $M$

% manuscript position: 26
c$'$) $F$ est $\pi$-accessible dans $\widehat C$,
% manuscript position: foot
% archive sheet number: 190
